% ============================================================================
% APPENDIX DV: LIT-DELLAVIGNA - Contracts, Media & Field Evidence
% ============================================================================
% Category: Literature Integration (LIT-)
% Language: English
% EBF Integration: Behavioral contracts, media effects, and field experiments
% ============================================================================

\chapter*{Appendix DV: Stefano DellaVigna Research Integration}
\addcontentsline{toc}{chapter}{Appendix DV: LIT-DELLAVIGNA - Contracts, Media \& Field Evidence}
\label{app:dellavigna}

% ============================================================================
% HEADER BLOCK
% ============================================================================
\begin{tcolorbox}[colback=white,colframe=black!75!white,title=Appendix Metadata]
\textbf{Appendix:} DV -- LIT-DELLAVIGNA \\
\textbf{Category:} Literature Integration (LIT-) \\
\textbf{Version:} 1.0 (January 2026) \\
\textbf{Dependencies:} CORE-WHEN, LIT-LAIBSON, LIT-LIST, METHOD-E \\
\textbf{Status:} Complete
\end{tcolorbox}

% ============================================================================
% CROSS-REFERENCE MAP
% ============================================================================
\begin{tcolorbox}[colback=blue!5!white,colframe=blue!75!black,title=Cross-Reference Map]
\textbf{Builds On:}
\begin{itemize}
    \item Appendix LA (LIT-LAIBSON): Present bias and $\beta$-$\delta$ model
    \item Appendix AB (LIT-LIST): Field experiments methodology
    \item Appendix RA (LIT-RABIN): Reference-dependent preferences
    \item Appendix V (CORE-WHEN): Temporal context effects
\end{itemize}

\textbf{Extends To:}
\begin{itemize}
    \item Appendix AD (DOMAIN-FINANCE): Consumer financial decisions
    \item Appendix AE (DOMAIN-HEALTH): Health behavior contracts
    \item Appendix AI (CONTEXT-MEDIA): Media influence on behavior
    \item Appendix S (PREDICT-POLICY): Nudge unit effectiveness
\end{itemize}

\textbf{Methods Connection:}
\begin{itemize}
    \item Appendix E (METHOD-E): Field experimental methodology
    \item Appendix AN (METHOD-LLMMC): Structural estimation
\end{itemize}
\end{tcolorbox}

% ============================================================================
% CHAPTER LINKAGE
% ============================================================================
\begin{tcolorbox}[colback=green!5!white,colframe=green!75!black,title=Chapter Linkage]
\textbf{Primary:} Chapter 11 (Time Preferences), Chapter 14 (Policy Applications) \\
\textbf{Secondary:} Chapter 9 (Incentives), Chapter 10 (Context Effects)
\end{tcolorbox}

% ============================================================================
% ABSTRACT
% ============================================================================
\begin{abstract}
This appendix integrates Stefano DellaVigna's research on behavioral contracts, media effects, and field experimental evidence into the Evidence-Based Framework. DellaVigna's work bridges theory and empirical evidence, demonstrating how present-biased consumers make predictable mistakes in contract choice (gym memberships), how media systematically influences political behavior (Fox News effect), and establishing benchmarks for what behavioral interventions actually achieve at scale (nudge unit meta-analysis). This appendix provides critical empirical parameters for CORE-WHEN and establishes methodological standards for field evidence in behavioral economics.
\end{abstract}

% ============================================================================
% QUICK REFERENCE
% ============================================================================
\begin{tcolorbox}[colback=gray!10!white,colframe=gray!50!black,title=Quick Reference]
Key terms defined in this appendix (see also Appendix G for complete Glossary):
\begin{description}
    \item[Behavioral Contract] Contract exploiting or helping with consumer biases
    \item[Overconfidence in Self-Control] Underestimating future self-control problems
    \item[Media Persuasion Effect] Causal impact of media exposure on attitudes/behavior
    \item[Nudge Unit] Government behavioral insights team implementing RCTs at scale
\end{description}
\end{tcolorbox}

% ============================================================================
% FUNDAMENTAL QUESTION
% ============================================================================
% -----------------------------------------------------------------------------
% APPENDIX SCOPE BOX
% -----------------------------------------------------------------------------

\begin{tcolorbox}[
    colback=orange!5!white,
    colframe=orange!75!black,
    title={\textbf{Appendix Scope: Ziel / In-Scope / Out-of-Scope / Constraints}},
    fonttitle=\bfseries
]

\textbf{Ziel (Objective):}
\begin{itemize}[nosep]
    \item How do systematic behavioral biases manifest in real-world markets and policy contexts, and what does field evidence tell us about their magnitude and malleability?
\end{itemize}

\textbf{In-Scope:}
\begin{itemize}[nosep]
    \item Fundamental Question
    \item Theoretical Framework
    \item Core Axioms
    \item Key Empirical Results
    \item EBF Integration
\end{itemize}

\textbf{Out-of-Scope (delegiert an andere Appendices):}
\begin{itemize}[nosep]
    \item LIT-LAIBSON $\rightarrow$ Appendix LA
    \item LIT-LIST $\rightarrow$ Appendix AB
    \item LIT-RABIN $\rightarrow$ Appendix RA
    \item CORE-WHEN $\rightarrow$ Appendix V
    \item DOMAIN-FINANCE $\rightarrow$ Appendix AD
\end{itemize}

\textbf{Constraints (Anwendungsgrenzen):}
\begin{itemize}[nosep]
    \item [TODO: Define when this appendix does NOT apply]
    \item [TODO: Define required prerequisites]
    \item [TODO: Define known limitations]
\end{itemize}

\textbf{Lieferobjekte (Deliverables):}
\begin{itemize}[nosep]
    \item 5 Table(s)
    \item 11 Equation(s)
    \item 6 Axiom(s)
    \item Worked Example(s)
\end{itemize}

\end{tcolorbox}


\section{Fundamental Question}

\begin{tcolorbox}[colback=yellow!10!white,colframe=orange!75!black,title=Central Research Question]
\textbf{How do systematic behavioral biases manifest in real-world markets and policy contexts, and what does field evidence tell us about their magnitude and malleability?}
\end{tcolorbox}

This question bridges theory (CORE-WHEN) and evidence (METHOD-E).

% ============================================================================
% THEORETICAL FRAMEWORK
% ============================================================================
\section{Theoretical Framework}

\subsection{Behavioral Contract Design}

DellaVigna \& Malmendier model contract choice under present bias and overconfidence:

\begin{equation}
\text{Contract Choice} = f(\beta, \hat{\beta}, \text{Usage Forecast}, \text{Price Structure})
\end{equation}

Where:
\begin{itemize}
    \item $\beta$ = True present bias parameter
    \item $\hat{\beta}$ = Perceived present bias (typically $\hat{\beta} > \beta$)
    \item Overconfidence: $\hat{\beta} - \beta > 0$ (think future self has more self-control)
\end{itemize}

\subsection{Media Persuasion Model}

\begin{equation}
\text{Vote}_i = \alpha + \gamma \cdot \text{Fox Exposure}_i + X_i'\beta + \epsilon_i
\end{equation}

Identification strategy: Exploit geographic variation in Fox News channel availability.

\subsection{Effort and Motivation}

From DellaVigna \& Pope's large-scale experiment:

\begin{equation}
\text{Effort} = f(\text{Intrinsic}, \text{Monetary}, \text{Social}, \text{Psychological})
\end{equation}

Different motivators have heterogeneous effects that experts systematically mispredict.

% ============================================================================
% AXIOMS
% ============================================================================
\section{Core Axioms}

\begin{tcolorbox}[colback=red!5!white,colframe=red!75!black,title=Axiom DV-1: Overconfidence in Self-Control]
\textbf{Consumers systematically overestimate their future gym attendance and underestimate their susceptibility to present bias.}

\begin{equation}
\frac{\text{Predicted Attendance}}{\text{Actual Attendance}} \approx 2.3
\end{equation}

Monthly members pay \$17/visit vs. \$10/visit pay-per-visit option.

\textit{Evidence:} DellaVigna \& Malmendier (2006), AER, $N=7,752$ members.
\end{tcolorbox}

\begin{tcolorbox}[colback=red!5!white,colframe=red!75!black,title=Axiom DV-2: Media Persuasion Effect]
\textbf{Media exposure causally affects political behavior.}

\begin{equation}
\frac{\partial \text{Republican Vote Share}}{\partial \text{Fox News}} \approx 0.4-0.7 \text{ pp}
\end{equation}

Fox News introduction increased Republican vote share by 0.4-0.7 percentage points.

\textit{Evidence:} DellaVigna \& Kaplan (2007), QJE.
\end{tcolorbox}

\begin{tcolorbox}[colback=red!5!white,colframe=red!75!black,title=Axiom DV-3: Social Pressure Avoidance]
\textbf{People avoid situations where they would face social pressure to give.}

\begin{equation}
P(\text{open door}) | \text{charity solicitor expected} < P(\text{open door}) | \text{no solicitor}
\end{equation}

Door-opening decreased 10-25\% when charity solicitation was expected.

\textit{Evidence:} DellaVigna, List \& Malmendier (2012), QJE.
\end{tcolorbox}

\begin{tcolorbox}[colback=red!5!white,colframe=red!75!black,title=Axiom DV-4: Motivation Heterogeneity]
\textbf{Different motivational treatments have highly heterogeneous effects on effort.}

\begin{equation}
\sigma^2_{\text{treatment effects}} >> \sigma^2_{\text{expert predictions}}
\end{equation}

Monetary incentives, social comparisons, and psychological framings differ more than experts predict.

\textit{Evidence:} DellaVigna \& Pope (2018), REStud, $N=9,861$ workers.
\end{tcolorbox}

\begin{tcolorbox}[colback=red!5!white,colframe=red!75!black,title=Axiom DV-5: Reference-Dependent Job Search]
\textbf{Unemployment benefits create reference points that affect job search behavior.}

\begin{equation}
\text{Hazard Rate}(t) \uparrow \text{ sharply at benefit exhaustion}
\end{equation}

Job finding spikes at benefit exhaustion, consistent with reference-dependent utility.

\textit{Evidence:} DellaVigna et al. (2017), QJE, Hungarian administrative data.
\end{tcolorbox}

\begin{tcolorbox}[colback=red!5!white,colframe=red!75!black,title=Axiom DV-6: Nudge Attenuation at Scale]
\textbf{Behavioral interventions show smaller effects when scaled from academic trials to government implementation.}

\begin{equation}
\frac{\text{Effect}_{nudge unit}}{\text{Effect}_{academic}} \approx 0.24
\end{equation}

Average nudge unit effect: 1.4 pp vs. academic average: 8.7 pp.

\textit{Evidence:} DellaVigna \& Linos (2022), Econometrica, 126 RCTs.
\end{tcolorbox}

% ============================================================================
% KEY EMPIRICAL RESULTS
% ============================================================================
\section{Key Empirical Results}

\subsection{The Gym Membership Study}

\begin{table}[h]
\centering
\begin{tabular}{lccc}
\hline
\textbf{Contract Type} & \textbf{Monthly Fee} & \textbf{Avg Visits} & \textbf{Cost/Visit} \\
\hline
Monthly (auto-renew) & \$70 & 4.3 & \$17.27 \\
Annual (prepaid) & \$58/mo & 4.8 & \$14.50 \\
Pay-per-visit & -- & -- & \$10.00 \\
\hline
\end{tabular}
\caption{Gym contract choices reveal overconfidence}
\end{table}

\textbf{Key insight:} 80\% of monthly members would save money with pay-per-visit.

\subsection{What Motivates Effort?}

\begin{table}[h]
\centering
\begin{tabular}{lcc}
\hline
\textbf{Treatment} & \textbf{Effect (buttons)} & \textbf{Expert Prediction} \\
\hline
4 cents/button (baseline) & 1,521 & 1,488 \\
Social comparison & +149 & +62 \\
Charity (1 cent/button) & +98 & +120 \\
``You are doing great'' & +40 & +24 \\
Task significance & +27 & +74 \\
\hline
\end{tabular}
\caption{Actual vs. predicted motivational effects}
\end{table}

\textbf{Key insight:} Experts underpredict social comparison, overpredict task significance.

\subsection{Nudge Units at Scale}

\begin{table}[h]
\centering
\begin{tabular}{lccc}
\hline
\textbf{Source} & \textbf{N trials} & \textbf{Avg Effect} & \textbf{Publication Bias?} \\
\hline
Academic papers & 74 & 8.7 pp & Yes (likely) \\
US Nudge Unit (OES) & 67 & 1.4 pp & No \\
UK Nudge Unit (BIT) & 59 & 2.0 pp & Minimal \\
\hline
\end{tabular}
\caption{Nudge effects by source}
\end{table}

\textbf{Key insight:} Publication bias and scaling challenges explain the gap.

% ============================================================================
% EBF INTEGRATION
% ============================================================================
\section{EBF Integration}

\subsection{CORE-WHEN Parameter Refinement}

DellaVigna's work refines temporal parameters:

\begin{equation}
\Psi_t^{DV} = \begin{pmatrix}
\hat{\beta} - \beta & \text{(overconfidence gap)} \\
\gamma_{media} & \text{(media persuasion)} \\
\alpha_{attenuation} & \text{(scale attenuation)} \\
\end{pmatrix}
\end{equation}

\begin{table}[h]
\centering
\begin{tabular}{lccl}
\hline
\textbf{Parameter} & \textbf{Estimate} & \textbf{SE} & \textbf{Source} \\
\hline
$\hat{\beta} - \beta$ & 0.30 & 0.05 & Gym study (2006) \\
$\gamma_{Fox}$ & 0.006 & 0.002 & Fox News (2007) \\
$\alpha_{scale}$ & 0.24 & 0.08 & Nudge units (2022) \\
Social pressure avoid & 0.15 & 0.04 & Charity (2012) \\
\hline
\end{tabular}
\caption{EBF parameters from DellaVigna research}
\end{table}

\subsection{10C Framework Mapping}

\begin{table}[h]
\centering
\begin{tabular}{lll}
\hline
\textbf{CORE} & \textbf{DellaVigna Contribution} & \textbf{Key Paper} \\
\hline
WHO (AAA) & Social pressure on giving & QJE 2012 \\
WHAT (C) & Reference-dependent job search & QJE 2017 \\
HOW (B) & Motivation heterogeneity & REStud 2018 \\
WHEN (V) & Overconfidence in $\hat{\beta}$ & AER 2006 \\
WHERE (BBB) & Scale attenuation factor & Econometrica 2022 \\
AWARE (AU) & Naivete about self-control & QJE 2004 \\
\hline
\end{tabular}
\caption{10C Framework mapping for DellaVigna research}
\end{table}

% ============================================================================
% WORKED EXAMPLE
% ============================================================================
\section{Worked Example: Predicting Intervention Effects at Scale}

\begin{tcolorbox}[colback=blue!5!white,colframe=blue!75!black,title=Application: Scaling a Retirement Savings Nudge]
\textbf{Context:} A government agency wants to scale an academic nudge showing +15pp enrollment.

\textbf{Step 1: Apply Attenuation Factor}

Using DellaVigna \& Linos (2022):
\begin{equation}
\text{Expected Effect}_{scale} = 15pp \times 0.24 = 3.6pp
\end{equation}

\textbf{Step 2: Adjust for Context}

Consider additional factors:
\begin{itemize}
    \item Academic sample: Self-selected, high attention
    \item Government sample: Universal, low attention
    \item Implementation fidelity: Typically 70-80\%
\end{itemize}

Revised estimate: $3.6pp \times 0.75 = 2.7pp$

\textbf{Step 3: Cost-Benefit Analysis}

\begin{align}
\text{Per-person cost} &= \$2 \text{ (mailing + processing)} \\
\text{Per-enrollment value} &= \$500 \text{ (lifetime tax revenue)} \\
\text{Break-even conversion} &= \frac{\$2}{\$500} = 0.4pp
\end{align}

Since $2.7pp > 0.4pp$, intervention is cost-effective.

\textbf{Step 4: Set Realistic Expectations}

Report to stakeholders:
\begin{itemize}
    \item Academic benchmark: 15pp (likely inflated)
    \item Realistic range: 2-5pp
    \item Still highly cost-effective
\end{itemize}
\end{tcolorbox}

% ============================================================================
% IMPLICATIONS
% ============================================================================
\section{Implications for Practice}

\subsection{For Contract Design}

\begin{enumerate}
    \item \textbf{Anticipate Overconfidence:} Consumers will overestimate future usage
    \item \textbf{Offer Commitment:} Allow opt-in to restriction contracts
    \item \textbf{Default to Protection:} Auto-cancel unused subscriptions
    \item \textbf{Transparent Pricing:} Show per-use costs clearly
\end{enumerate}

\subsection{For Policy Scaling}

\begin{enumerate}
    \item \textbf{Expect Attenuation:} Academic effects shrink $\sim$75\% at scale
    \item \textbf{Pre-Register Trials:} Reduce publication bias
    \item \textbf{Invest in Fidelity:} Implementation quality matters
    \item \textbf{Multiple Replications:} Don't rely on single studies
\end{enumerate}

% ============================================================================
% SUMMARY
% ============================================================================
\section{Summary}

\begin{tcolorbox}[colback=gray!10!white,colframe=gray!50!black,title=Key Takeaways]
\begin{enumerate}
    \item \textbf{Overconfidence Gap:} $\hat{\beta} - \beta \approx 0.30$ (think 2.3x more visits)
    \item \textbf{Media Effects:} Fox News $\rightarrow$ +0.5pp Republican vote
    \item \textbf{Social Pressure:} People actively avoid giving situations
    \item \textbf{Scale Attenuation:} Nudge effects shrink 76\% from lab to field
    \item \textbf{Motivation Variance:} Experts mispredict heterogeneity
    \item \textbf{Reference Points:} Job search spikes at benefit exhaustion
\end{enumerate}

\textbf{Core Insight:} Field evidence reveals systematic biases are real but interventions attenuate at scale.
\end{tcolorbox}

% ============================================================================
% GLOSSARY
% ============================================================================
\section{Glossary}

Key terms defined in this appendix (see also Appendix G for complete Glossary):

\begin{description}
    \item[Behavioral Contract] Contract designed around consumer biases
    \item[Naivete] Unawareness of own present bias ($\hat{\beta} > \beta$)
    \item[Sophistication] Awareness of own present bias ($\hat{\beta} = \beta$)
    \item[Partial Naivete] Underestimating present bias ($\beta < \hat{\beta} < 1$)
    \item[Scale Attenuation] Reduction in effect sizes when scaling interventions
    \item[Nudge Unit] Government team running behavioral RCTs
    \item[Social Pressure] Utility cost of refusing prosocial requests
\end{description}

% ============================================================================
% CRITICAL FOUNDATIONS
% ============================================================================
\section{Critical Foundations}

\begin{enumerate}
    \item \textbf{Laibson (1997):} $\beta$-$\delta$ model of hyperbolic discounting
    \item \textbf{O'Donoghue \& Rabin (2001):} Naivete and sophistication
    \item \textbf{Kőszegi \& Rabin (2006):} Reference-dependent preferences
    \item \textbf{List (2011):} Why field experiments matter
    \item \textbf{Thaler \& Sunstein (2008):} Nudge and choice architecture
\end{enumerate}

% ============================================================================
% OPEN ISSUES
% ============================================================================
\section{Open Issues}

\begin{enumerate}
    \item \textbf{Attenuation Mechanism:} What exactly causes scale attenuation?
    \item \textbf{Heterogeneity:} Which populations benefit most from nudges?
    \item \textbf{Long-term Effects:} Do behavioral interventions persist?
    \item \textbf{Welfare:} Are nudges paternalistic or welfare-enhancing?
    \item \textbf{Media Evolution:} How do social media effects compare to TV?
\end{enumerate}

% ============================================================================
% REFERENCES
% ============================================================================
\section{References}

See Master Bibliography (bcm2\_master\_references.bib) for complete citations.

\nocite{bcm_master}

\textbf{Primary DellaVigna Sources} (23 papers integrated):

\begin{itemize}
    \item DellaVigna, S. \& Malmendier, U. (2004). Contract Design and Self-Control: Theory and Evidence. \textit{QJE}, 119(2). [NTM: Time inconsistency exploitation]
    \item DellaVigna, S. \& Malmendier, U. (2006). Paying Not to Go to the Gym. \textit{AER}, 96(3).
    \item DellaVigna, S. \& Kaplan, E. (2007). The Fox News Effect. \textit{QJE}, 122(3).
    \item DellaVigna, S. (2009). Psychology and Economics: Evidence from the Field. \textit{JEL}, 47(2).
    \item DellaVigna, S., List, J. \& Malmendier, U. (2012). Testing for Altruism and Social Pressure. \textit{QJE}, 127(1).
    \item DellaVigna, S. et al. (2017). Reference-Dependent Job Search. \textit{QJE}, 132(4).
    \item DellaVigna, S. \& Pope, D. (2018). What Motivates Effort? \textit{REStud}, 85(2).
    \item DellaVigna, S. \& Linos, E. (2022). RCTs to Scale. \textit{Econometrica}, 90(1).
    \item DellaVigna, S. (2020). Structural Behavioral Economics. \textit{Handbook of Behavioral Economics}.
\end{itemize}

\end{document}
